%%%%%%%%%%%%%%%%%%%%%%%%%%%%%%%%%%%%%%%%%%%%%%%%%%%%%%%%%%%%%%%%%%%%%%%%%%%
\chapter{Review of Related Systems and Literature}
%%%%%%%%%%%%%%%%%%%%%%%%%%%%%%%%%%%%%%%%%%%%%%%%%%%%%%%%%%%%%%%%%%%%%%%%%%%

This chapter briefly discusses various scientific journals and systems that are related to the benefits of astronomy and space 
exploration, the advantages and disadvantages of game-based education, and the effects of collaborative learning. Research was
 collected and thoroughly analyzed to determine the problem gap that Aphelion will attempt to fill, as well as to synthesize 
 insights to create an optimal balance between fun and learning. Optimization methods and graph theory are briefly discussed 
 to calculate the advancement between stars and to ensure a smooth experience in the game. Additionally software that are 
 related to the goal of Aphelion have been investigated, to further determine if the perceived problem gap is viable. The 
 related systems run in platforms such as mobile and personal computer, and are available in the Android Play Store, Flathub 
 for Linux machines, or as standalone .exe files.


\section{Research-Related Literature}

\subsection{Instructional Design Frameworks}

Instructional design frameworks provide a systemic approach to developing learning tools for students that are 
integrated with technological features and digital interfaces [22]. The ADDIE model is a popular example of an 
instructional design framework that consists of Analysis, Design, Development, Implementation, and Evaluation [23][28]. 
Each stage encourages iteration and continuous improvement. The Analysis stage identifies the assets in a learning 
environment – requirements, student needs, objectives, and learning materials are examples of variables that will be 
identified in this stage. Next is the Design stage, where storyboards are brainstormed and media formats and templates 
are agreed upon using the results of the previous stage as a context. Development is when these abstract concepts are 
turned to concrete products using a combination of software and talent by developers, then assessed for quality control. 
The system is then released into production, where it will be used by classroom and online platforms, which is the 
Implementation stage, and subsequently the Evaluation stage where the system’s effectiveness in learning, and other 
feedback are quantified [22].

\subsection{Gamification and Game-Based Learning (GBL) Strategies}

Gamification and Game-Based Learning both attempt to enhance user engagement and motivation in learning systems [18][20][25]. 
However, both strategies differ in the implementation. Gamified systems focus on adding elements of game design in environments 
that traditionally do not have game context, such as in learning tools, for the purpose of increasing user motivation and interaction. 
GBL on the other hand designs a full feature game system for the purpose of integrating educational content in it.

A gamified system typically consists of standardized elements to achieve its purpose [11]. A point system connected to a leaderboard 
is one of the most common elements, highlighting students that are excelling in their activities. However, this element may also lead 
to underachievers being demotivated when they see themselves in the bottom of the leaderboard [26]. Additionally, social interaction 
is implemented to great effect – competition drives students to work harder to land at the top and be praised for their performance, 
while collaboration promotes a sense of community and can help develop social skills and critical thinking as a group, as problems are 
solved with team effort. Community-based learning platforms may offer increased effectiveness in studying, evidenced by the transformation
of a Brazilian educational platform from a centralized model to a decentralized, community-based system encouraged by gamified elements,
leading to a self-sufficient community network [12].  Progress indicators, such as leveling systems and progress bars are used to provide
a concrete and measurable sense of progress that easily reflects a student’s effort and competence. Lastly, adding a narrative storyline
helps the system to be immersive and allows students to roleplay, increasing the sense of belonging. [11]

A system based on GBL typically has most, if not all of the elements that gamified systems have. The learning process turns to a game 
environment with a feedback loop of students facing challenges, and having solutions be designed, either alone or through collaboration, 
applying what they learned in the process. [11]

In both systems, various social strategies pertain to the interaction of students. Through collaboration, students may exhibit a better 
learning performance and flow experience [11], while an element of competition may promote the drive of self-efficacy, focus,  and rapid 
improvement to meet the competitive pressure during the game [29]. Both types of interaction are generally considered to be superior to a 
solo learning environment in terms of educational performance and student motivation [30]. Educational games, especially modern ones, also
benefit from practicing adaptive gamification, where individual students can dynamically set the difficulty of the game to be easier or 
harder to match their skill [11]. In line with individualistic elements, allowing students to express themselves through avatars supports 
their autonomy [18].

A common problem observed in an Indonesian study is that despite the presence of high-tech tools for instructors to utilize, such as 
WiFi, and cheap computers for students such as Chromebooks, traditional printed materials are still the preferred tools, minimizing the 
benefits of the newer tools. Additionally, the style of traditional teaching does not blend well with the fast-paced and modern design of 
GBL systems [15]. By switching away from traditional teaching strategies, students may show more enthusiasm in the new way of learning 
that features more interaction and collaboration [21]. In fact, a survey from a research paper studying the effectiveness of game-based 
learning noted that 95\% of students who interacted with game-based systems found it engaging and boosted their motivation [14]. 
This shows that GBL may enhance higher-level thinking skills by immersing students in friendly play, especially if the students perceive 
the system as useful enough to learn or provides positive emotions as a reward [19]. Additionally, clear, visual-heavy interactions in 
games is reportedly a “fun experience”, said by students [9].

The efficacy of these strategies is rooted in three major theories. The Self-Determination Theory claims that the satisfaction of 
autonomy, competence, and relatedness, agreed to be part of the basic needs of humans, are satisfied. This leads to the Theory of 
Gamified Learning, which suggests that game elements, because of their effectiveness in satisfying basic human needs, encourage reward 
seeking in learning systems. Flow Theory is a complementary theory that describes a psychological state of players experiencing intense 
concentration and enjoyment when playing, giving the feeling that they are “one” with the game, often happening when goals and challenges 
are clear [18][27].

The concepts of gamification and game-based learning are important for Aphelion, because it may help students learn about astronomy in 
an immersive and fun way. By providing multiple avenues for competition, collaboration, self-expression, points, levels, and other game 
elements, interwoven with technology, Aphelion aims to be an educational platform with engineered game elements for students.

\subsection{Current State of Astronomy Education and Curriculum}

There are demands for better representation of astronomy in the curriculum, as there are gaps in the instruction of the subject 
[24]. Primary and middle school focus on celestial motion and the solar system, and then large-scale cosmology in upper high school,
which leaves the students with no conceptual understanding of the Milky Way, nearby star systems, or the immediate vicinity of galaxies 
beyond our own, as this structure overlooks the local galactic and extragalactic environments [24].

There is also a challenge in acquiring formal and credible knowledge, as some often derive from informal sources such as social media, 
movies, and unverified internet content [13]. These informal sources can spread speculative information, which undermines the work of 
formal educators. To address this problem, the literature advocates for authentic investigations that move beyond passive lectures [13]. 

\subsection{Digital Platforms and Virtual Environments in Science}

Integrating digital platforms and virtual environments into formal education expands learning opportunities through a dynamic 
and interactive process. Utilizing this technology is crucial in disciplines like astronomy, where being able to visualize and 
simulate is essential in grasping abstract concepts that cannot be demonstrated in traditional methods.

Virtual environments allow students to visualize and manipulate objects and concepts in an accessible way; rather than going to 
laboratories, they can use the natural physics in software like games to simulate science-related concepts for exploring planetary 
collisions and realistic rocketry [1][5][8][10]. They can also be used as observation tools with platforms such as Stellarium and 
Skysafari, which enable real-time tracking of astronomical objects. With these platforms, it is possible to increase student interest 
in learning about the universe, and it even gives them a deeper understanding of concepts [17].

Mobile gamified applications serve as a popular tool for primary education among Generation Alpha digital natives, as most of them 
have used mobile devices before they are one year old [11]. Interactivity is essential for learning among younger generations, as it 
increases their curiosity and knowledge retention; thus, the use of AR-based mobile applications is prevalently used to teach astronomy 
and plant studies [6] [11].

\section{Review of Related Systems}

\subsection{Feature-Complete, Free Open Source (FOSS) Planetarium}

Stellarium provides a 3D simulation of the night sky through a telescope. The default version already includes over 600,000 stars 
and 80,000 deep-sky objects, such as nebulae, with each celestial object being able to be explored and its information seen by clicking 
in the night sky. It can be expanded to over 210 million stars and 1 million deep-sky objects. The application also offers a full 
location and time control, allowing users to choose the night sky anywhere in the world, at any time of the day; in fact, the location 
and time is not limited to planet Earth – users can go to moons and other planets – some even being outside of the solar system – to 
stare at the planetary night sky. It also features a realistic atmosphere, with sunrises and sunsets, twinkling stars, shooting stars, 
and other atmospheric phenomena, as well as constellations and asterisms for 40+ different cultures. All of these features are highly 
customizable and can be enabled and disabled to one’s liking [4].

Another FOSS planetarium application is KStars. It focuses more on the needs of astrophotographers, with features such as an observation 
planner, graphical simulation, and astrophotography workflows. Its stars number up to 100 million, compared to Stellarium’s 600,000 to 
210 million stars. While Stellarium is the go-to application for a general study of the night sky, KStar finds its niche for more 
professional applications. [2]

\subsection{Space Combat Game With Star System Exploration}

Starsector is an open-world, space roleplaying game that primarily focuses on theatrical space combat. However, it features an in-depth 
exploration system, where the player’s fleet can fly into star systems, and scan individual planets to reveal their attributes. Some 
planets are habitable, some are too hot or cold, some have crushing gravity or toxic plants, all contributing to a planet’s hazard 
rating. Additionally, players can encounter anomalies in a star system, such as asteroid fields or belts. Most star systems are simulated 
models that are procedurally generated, sometimes being red dwarves, red giants, yellow stars, neutron stars, or even black holes. There 
is a sense of exploration as users visit a star system, scan its planets, and note for anomalies in space. [3]

\subsection{Physics-Based Space Simulator}

Universe Sandbox allows users to interact with celestial bodies – being able to create and destroy planets in a simulated space, 
using N-body simulation at any temporal speed. Users can build their own solar systems, complete with objects from moons to black holes, 
and with the ability to manipulate time, enable them to see the birth and death of stars. Planets can be made out of many materials such 
as hydrogen for gas planets, iron, rock, water, and more. Additionally, the climate of Earth is simulated and can be changed by changing 
its axial tilt or position relative to the sun. The application is a useful tool for simulating universal phenomena, but is notoriously 
difficult to learn, and lacks objectives for learning astronomy. The game prioritizes observation and visualization of celestial bodies 
more than actual gameplay. [5]

\subsection{Space Program Management Simulator}

Kerbal Space Program may be the most low-level and practical simulation of space travel available. It is a space program management 
simulator that lets players design and launch spacecraft with realistic aerodynamic and orbital physics. Players fully manage a space 
program to explore the game’s solar system, with budget mechanics implemented. The game dives into concepts of orbital dynamics and 
rocket science, allowing players to actually simulate two spacecraft docking, or landing on other planets. [1]

\subsection{Matrix Comparison}

\begin{table}[h!]
\centering
\setlength{\tabcolsep}{5pt}
\renewcommand{\arraystretch}{5}
\resizebox{\textwidth}{!}{%
\begin{tabular}{|c|c|c|c|c|c|c|}
\hline
\textbf{System} & \textbf{Aphelion} & \textbf{Stellarium} & \textbf{KStars} & \textbf{Starsector} & \textbf{Universe Sandbox} & \textbf{Kerbal Space Program}\\
\hline
Information about celestial bodies & $\checkmark$ & $\checkmark$ & $\checkmark$ & $\checkmark$ & $\checkmark$ & $\checkmark$\\
\hline
In-game collaboration & $\checkmark$ & $\times$ & $\times$ & $\times$ & $\times$ & $\times$\\
\hline
Objectives & $\checkmark$ & $\times$ & $\times$ & $\checkmark$ & $\times$ & $\checkmark$\\
\hline
Accessible learning curve & $\checkmark$ & $\times$ & $\times$ & $\times$ & $\times$ & $\times$\\
\hline
Portable & $\checkmark$ & $\checkmark$ & $\times$ & $\times$ & $\times$ & $\times$\\
\hline
Account creation and customization & $\checkmark$ & $\times$ & $\times$ & $\times$ & $\times$ & $\times$\\
\hline
\end{tabular}}
\caption{Comparison of Aphelion with related systems}
\end{table}